\documentclass[a4paper,12pt]{article}

\usepackage[T2A]{fontenc}
\usepackage[utf8]{inputenc}
\usepackage[russian]{babel}
\usepackage{amsmath,amssymb,amsthm,mathtools}
\usepackage{helvet}
\renewcommand{\familydefault}{\sfdefault}
\usepackage{icomma}
\usepackage[a4paper,margin=2cm,headheight=25pt]{geometry}
\usepackage{microtype}
\usepackage{tikz}
\usepackage{tkz-euclide}
\usepackage{pgfplots}
\pgfplotsset{compat=1.18}
\usepackage{tabularx,booktabs,longtable}
\usepackage{enumitem}
\usepackage{hyperref}
\usepackage{parskip}
\usepackage{fancyhdr}
\usepackage{setspace}
\usepackage{xcolor}
\usepackage{titlesec}

\definecolor{accent}{HTML}{008080}
\definecolor{darkaccent}{HTML}{004D4D}
\definecolor{textgray}{HTML}{333333}
\definecolor{softgray}{HTML}{EEF3F3}
\definecolor{warn}{HTML}{7A4B32}

\hypersetup{
  colorlinks=true,
  linkcolor=darkaccent,
  urlcolor=darkaccent,
  pdfauthor={Лёвин Артём Александрович},
  pdftitle={Чек-лист: текстовые задачи на движение}
}

\onehalfspacing
\setlength{\parindent}{0pt}
\setlength{\emergencystretch}{3em}
\setlist[itemize]{leftmargin=1.6em,itemsep=0.25em,topsep=0.35em}
\setlist[enumerate]{leftmargin=1.8em,itemsep=0.45em,topsep=0.35em}
\renewcommand{\arraystretch}{1.3}

\titleformat{\section}
  {\Large\bfseries\color{darkaccent}}
  {\thesection.}{0.6em}{}
\titleformat{\subsection}
  {\large\bfseries\color{accent}}
  {\thesubsection.}{0.55em}{}
\titlespacing*{\section}{0pt}{1.2em}{0.55em}
\titlespacing*{\subsection}{0pt}{0.9em}{0.35em}

\pagestyle{fancy}
\fancyhf{}
\fancyhead[L]{\small\color{textgray}Текстовые задачи на движение}
\fancyhead[R]{\small\color{textgray}Ксения Клыкова}
\fancyfoot[C]{\small\color{textgray}\thepage}
\renewcommand{\headrulewidth}{0.4pt}
\renewcommand{\headrule}{%
  \hbox to\headwidth{%
    \color{accent!55}\leaders\hrule height \headrulewidth\hfill
  }%
}

\newcommand{\unit}[1]{\,\text{#1}}

\newcommand{\formula}[1]{%
  \begin{center}
    \large\color{darkaccent}\ensuremath{#1}
  \end{center}
}

\newcommand{\checkpoint}[1]{%
  \par\smallskip
  {\color{warn}\textbf{Контроль:} #1}\par\smallskip
}

\begin{document}

%------------------------------------------------
% Титульный лист
%------------------------------------------------
\begin{titlepage}
  \thispagestyle{empty}
  \centering
  \vspace*{2.8cm}

  {\Huge\bfseries\color{darkaccent}Чек-лист\par}
  \vspace{0.6cm}

  {\Huge\bfseries Текстовые задачи на движение\par}
  \vspace{0.45cm}

  {\Large\color{textgray}
    ЕГЭ: модели, уравнения, единицы измерения\par
  }

  \vfill

  {\large Лёвин Артём Александрович\par}
  \vspace{0.2cm}

  {\large\bfseries эксклюзивно для Ксении Клыковой\par}
  \vspace{0.8cm}

  {\large \today\par}

  \begin{tikzpicture}[remember picture,overlay]
    \draw[accent!55,line width=1.2pt]
      ([xshift=1.5cm,yshift=1.35cm]current page.south west)
      .. controls
      ([xshift=5cm,yshift=2.05cm]current page.south west)
      and
      ([xshift=-5cm,yshift=0.75cm]current page.south east)
      ..
      ([xshift=-1.5cm,yshift=1.35cm]current page.south east);
  \end{tikzpicture}
\end{titlepage}

%------------------------------------------------
\section{Базовая модель движения}

Для равномерного движения используются три величины:
\formula{
  S=vt,
  \qquad
  v=\frac{S}{t},
  \qquad
  t=\frac{S}{v}.
}

\begin{table}[ht]
\centering
\caption{Смысл величин в задачах на движение}
\begin{tabularx}{\linewidth}{@{}l l X@{}}
\toprule
Величина & Обозначение & Смысл и типичные единицы \\
\midrule
Путь & $S$ & км, м; длина пройденного участка \\
Скорость & $v$ & км/ч, м/с; путь за единицу времени \\
Время & $t$ & ч, мин, с; продолжительность движения \\
\bottomrule
\end{tabularx}
\end{table}

\subsection{Универсальный алгоритм}

\begin{enumerate}
  \item Определите участников движения и отдельные этапы маршрута.
  \item Выберите общую систему единиц. Ориентиром служат единицы скорости.
  \item Введите переменную. По возможности обозначьте через $x$ искомую величину.
  \item Составьте таблицу $v$--$t$--$S$.
  \item Заполните два столбца данными, третий --- формулой движения.
  \item Найдите в условии связь между временами, путями или скоростями.
  \item Составьте и решите уравнение.
  \item Проверьте смысл ответа, единицы и требуемую величину.
\end{enumerate}

\checkpoint{
  Скорость положительна, время положительно, путь неотрицателен.
  Для движения против течения требуется $v_0-u>0$.
}

\subsection{Перевод единиц}

\formula{
  1\unit{ч}=60\unit{мин},
  \qquad
  1\unit{км}=1000\unit{м},
  \qquad
  1\unit{м/с}=3{,}6\unit{км/ч}.
}

Примеры:
\[
35\unit{мин}
=
\frac{35}{60}\unit{ч}
=
\frac{7}{12}\unit{ч},
\qquad
300\unit{м}
=
0{,}3\unit{км}.
\]

%------------------------------------------------
\section{Встречное движение и движение вдогонку}

\subsection{Относительная скорость}

\begin{center}
\begin{tikzpicture}[x=1cm,y=1cm,>=stealth]
  \draw[accent,line width=1pt,->]
    (0,1.2) -- (4,1.2)
    node[midway,above] {$v_1$};

  \draw[accent,line width=1pt,<-]
    (7,1.2) -- (11,1.2)
    node[midway,above] {$v_2$};

  \node[textgray] at (5.5,1.2)
    {навстречу: $v_{\text{отн}}=v_1+v_2$};

  \draw[accent,line width=1pt,->]
    (0,0) -- (4,0)
    node[midway,above] {$v_1$};

  \draw[darkaccent,line width=1pt,->]
    (0,-0.5) -- (2.7,-0.5)
    node[midway,below] {$v_2$};

  \node[textgray] at (7.1,-0.2)
    {одно направление: $v_{\text{отн}}=|v_1-v_2|$};
\end{tikzpicture}
\end{center}

При старте в разные моменты времени удобно сравнивать времена участников.
Если первый участник начал движение на $\tau$ часов раньше, то
\formula{
  t_1-t_2=\tau.
}

\subsection{Пример: встреча при запоздалом старте}

Расстояние между городами равно $435\unit{км}$.
Первый автомобиль выехал из города $A$ со скоростью
$60\unit{км/ч}$.
Через час навстречу ему из города $B$ выехал второй автомобиль
со скоростью $65\unit{км/ч}$.
Найдём расстояние от города $A$ до места встречи.

Пусть первый автомобиль прошёл $x\unit{км}$.
Тогда второй автомобиль прошёл $(435-x)\unit{км}$.

\begin{table}[ht]
\centering
\caption{Модель задачи о встрече}
\begin{tabularx}{\linewidth}{@{}X c c c@{}}
\toprule
Участник & $v$ & $t$ & $S$ \\
\midrule
Первый автомобиль
& $60$
& $\dfrac{x}{60}$
& $x$ \\
Второй автомобиль
& $65$
& $\dfrac{435-x}{65}$
& $435-x$ \\
\bottomrule
\end{tabularx}
\end{table}

Первый автомобиль находился в пути на один час дольше:
\[
\frac{x}{60}
-
\frac{435-x}{65}
=
1.
\]

После умножения на $60\cdot65$:
\[
65x-60(435-x)=3900,
\]
\[
65x-26100+60x=3900,
\]
\[
125x=30000,
\qquad
x=240.
\]

Ответ: $240\unit{км}$ от города $A$.

\subsection{Пример: увеличение расстояния между пешеходами}

Два пешехода одновременно вышли из одной точки в одном направлении.
Скорость первого на $1{,}5\unit{км/ч}$ больше скорости второго.
Через какое время расстояние между ними станет равным $300\unit{м}$?

Пусть время движения равно $x$ часов,
скорость медленного пешехода равна $y\unit{км/ч}$.
Тогда
\[
(y+1{,}5)x-yx=0{,}3,
\]
\[
1{,}5x=0{,}3,
\qquad
x=0{,}2\unit{ч}.
\]

Переведём часы в минуты:
\[
0{,}2\cdot60=12\unit{мин}.
\]

Ответ: $12\unit{мин}$.

\checkpoint{
  В задачах вдогонку расстояние между участниками изменяется
  со скоростью, равной разности их скоростей.
}

%------------------------------------------------
\section{Движение по круговой трассе}

Если два участника стартуют одновременно из одной точки и движутся
в одном направлении, первый полный обгон происходит тогда,
когда быстрый участник выигрывает ровно один круг:
\formula{
  \bigl(v_{\text{быстр}}-v_{\text{медл}}\bigr)t=L,
}
где $L$ --- длина круга.

После $k$ полных обгонов выигранный путь равен $kL$:
\[
\bigl(v_{\text{быстр}}-v_{\text{медл}}\bigr)t=kL.
\]

\subsection{Пример}

Длина трассы равна $21\unit{км}$.
Скорость первого автомобиля равна $54\unit{км/ч}$.
Через $35\unit{мин}$ второй автомобиль впервые обогнал первый.

Переведём время в часы:
\[
t
=
\frac{35}{60}
=
\frac{7}{12}\unit{ч}.
\]

Пусть скорость второго автомобиля равна $x\unit{км/ч}$.
Тогда
\[
(x-54)\cdot\frac{7}{12}=21,
\]
\[
x-54=21\cdot\frac{12}{7}=36,
\]
\[
x=90.
\]

Ответ: $90\unit{км/ч}$.

%------------------------------------------------
\section{Движение по воде}

Пусть $v_0$ --- собственная скорость судна,
$u$ --- скорость течения.

\formula{
  v_{\text{по течению}}=v_0+u,
  \qquad
  v_{\text{против течения}}=v_0-u.
}

Условие применимости:
\[
v_0>u.
\]

Стоянка входит в общее время путешествия.
Путь во время стоянки равен нулю.

\subsection{Пример}

Собственная скорость теплохода равна $25\unit{км/ч}$,
скорость течения --- $3\unit{км/ч}$.
Стоянка длится $5\unit{ч}$,
всё путешествие занимает $30\unit{ч}$.
Найдём путь за весь рейс.

Пусть расстояние в одну сторону равно $x\unit{км}$.

Скорость по течению:
\[
25+3=28\unit{км/ч}.
\]

Скорость против течения:
\[
25-3=22\unit{км/ч}.
\]

Уравнение:
\[
\frac{x}{28}
+
5
+
\frac{x}{22}
=
30.
\]

Перенесём время стоянки:
\[
\frac{x}{28}
+
\frac{x}{22}
=
25.
\]

Приведём дроби к общему знаменателю:
\[
x\left(
\frac{1}{28}
+
\frac{1}{22}
\right)
=
25,
\]
\[
x\left(
\frac{22+28}{616}
\right)
=
25,
\]
\[
x\cdot\frac{50}{616}
=
25,
\]
\[
x\cdot\frac{25}{308}
=
25,
\qquad
x=308.
\]

Расстояние за весь рейс:
\[
2x=616\unit{км}.
\]

Ответ: $616\unit{км}$.

\checkpoint{
  В ответе учитывайте точную формулировку:
  путь в одну сторону, обратный путь или весь рейс.
}

%------------------------------------------------
\section{Средняя скорость}

Средняя скорость на всём маршруте определяется формулой
\formula{
  v_{\text{ср}}
  =
  \frac{S_{\text{общ}}}{t_{\text{общ}}}.
}

Для нескольких участков:
\[
v_{\text{ср}}
=
\frac{S_1+S_2+\dots+S_n}
{\dfrac{S_1}{v_1}
+\dfrac{S_2}{v_2}
+\dots
+\dfrac{S_n}{v_n}}.
\]

\subsection{Равные участки пути}

Если участки пути равны, можно обозначить длину каждого участка через $x$.

Для трёх равных участков:
\formula{
  v_{\text{ср}}
  =
  \frac{3}
  {\dfrac{1}{v_1}
  +\dfrac{1}{v_2}
  +\dfrac{1}{v_3}}.
}

Для двух равных участков:
\[
v_{\text{ср}}
=
\frac{2v_1v_2}{v_1+v_2}.
\]

\subsection{Равные промежутки времени}

Если промежутки времени равны, средняя скорость равна
среднему арифметическому скоростей:
\[
v_{\text{ср}}
=
\frac{v_1+v_2+\dots+v_n}{n}.
\]

\subsection{Пример}

Первую треть пути велосипедист ехал со скоростью
$30\unit{км/ч}$,
вторую --- $20\unit{км/ч}$,
третью --- $15\unit{км/ч}$.

Пусть длина каждой трети равна $x$.
Тогда
\[
v_{\text{ср}}
=
\frac{3x}
{\dfrac{x}{30}
+\dfrac{x}{20}
+\dfrac{x}{15}}.
\]

Сократим на $x$:
\[
v_{\text{ср}}
=
\frac{3}
{\dfrac{1}{30}
+\dfrac{1}{20}
+\dfrac{1}{15}}.
\]

Общий знаменатель равен $60$:
\[
v_{\text{ср}}
=
\frac{3}
{\dfrac{2+3+4}{60}}
=
\frac{3}{\dfrac{9}{60}}
=
\frac{3\cdot60}{9}
=
20\unit{км/ч}.
\]

Ответ: $20\unit{км/ч}$.

\checkpoint{
  Среднее арифметическое скоростей используется при равных промежутках времени.
  Для равных участков пути вычисляйте отношение общего пути к общему времени.
}

%------------------------------------------------
\section{Движение протяжённых тел}

К протяжённым телам относятся поезда, колонны, суда и другие объекты,
длину которых требуется учитывать.

\subsection{Прохождение неподвижного объекта}

Если поезд длиной $l$ проходит мимо столба или наблюдателя, требуется,
чтобы весь поезд переместился на собственную длину:
\[
t=\frac{l}{v}.
\]

Если поезд проходит мост или тоннель длиной $L$, голова поезда должна пройти
расстояние $l+L$:
\[
t=\frac{l+L}{v}.
\]

\subsection{Два движущихся объекта}

При полном прохождении двух протяжённых объектов друг мимо друга
учитывается сумма их длин:
\[
t
=
\frac{l_1+l_2}{v_{\text{отн}}}.
\]

Если в условии дан дополнительный начальный промежуток $d$, получаем:
\[
t
=
\frac{l_1+l_2+d}{v_{\text{отн}}}.
\]

Для встречного движения:
\[
v_{\text{отн}}=v_1+v_2.
\]

Для движения в одном направлении:
\[
v_{\text{отн}}=|v_1-v_2|.
\]

\subsection{Пример с суммарным относительным путём}

Для завершения манёвра требуется компенсировать:

\begin{itemize}
  \item длину первого объекта $120\unit{м}$;
  \item длину второго объекта $80\unit{м}$;
  \item начальный промежуток $400\unit{м}$;
  \item дополнительный участок перегона $600\unit{м}$.
\end{itemize}

Суммарный относительный путь:
\[
120+80+400+600
=
1200\unit{м}
=
1{,}2\unit{км}.
\]

Время:
\[
12\unit{мин}
=
\frac{12}{60}\unit{ч}
=
0{,}2\unit{ч}.
\]

Разность скоростей:
\[
\Delta v
=
\frac{1{,}2}{0{,}2}
=
6\unit{км/ч}.
\]

Ответ: $6\unit{км/ч}$.

\checkpoint{
  Для протяжённых тел переведите длины и скорости
  в согласованные единицы до подстановки в формулу.
}

%------------------------------------------------
\section{Алгебраические приёмы}

\subsection{Работа с дробными уравнениями}

При решении уравнения
\[
\frac{x}{60}
-
\frac{435-x}{65}
=
1
\]
общий знаменатель равен $60\cdot65$.

После умножения всего уравнения на общий знаменатель:
\[
65x-60(435-x)=60\cdot65.
\]

Важно учитывать знак перед второй дробью:
\[
-60(435-x)
=
-26100+60x.
\]

\subsection{Рациональные вычисления}

\begin{itemize}
  \item Сохраняйте произведения в разложенном виде до этапа сокращения.
  \item Выносите общий множитель.
  \item Сокращайте числа до вычисления крупных произведений.
  \item После раскрытия скобок отдельно проверяйте каждый знак.
  \item После решения возвращайтесь к вопросу задачи.
\end{itemize}

Пример:
\[
7(x-54)=21\cdot12.
\]

Разделим обе части на $7$:
\[
x-54=3\cdot12=36.
\]

Такой порядок уменьшает объём арифметики.

%------------------------------------------------
\section{Типичные ошибки}

\begin{itemize}
  \item Минуты оставлены рядом со скоростью в километрах в час.
  \item Метры смешаны с километрами.
  \item При встречном движении использована разность скоростей.
  \item При движении вдогонку использована сумма скоростей.
  \item Время стоянки пропущено в общей продолжительности рейса.
  \item Скорости усреднены без проверки равенства времён.
  \item Найден путь в одну сторону, хотя требуется весь рейс.
  \item В задаче о круговой трассе потерян выигрыш одного полного круга.
  \item Длина поезда пропущена при прохождении моста.
  \item При полном обгоне пропущена длина одного из объектов.
  \item После раскрытия скобок потерян знак.
  \item Полученная величина записана в неверных единицах.
\end{itemize}

%------------------------------------------------
\section{Итоговый чек-лист}

Перед записью ответа проверьте:

\begin{enumerate}
  \item Все величины переведены в согласованные единицы.
  \item Таблица $v$--$t$--$S$ заполнена по формуле $S=vt$.
  \item Учтено направление движения.
  \item Учтена разница во времени старта.
  \item Для круговой трассы записано целое число выигранных кругов.
  \item Для движения по воде выполнено условие $v_0>u$.
  \item В средней скорости использованы общий путь и общее время.
  \item Для протяжённых тел учтены необходимые длины.
  \item Решение уравнения проверено арифметически.
  \item В ответе указана величина, требуемая условием.
\end{enumerate}

%------------------------------------------------
\clearpage
\section*{Тренировочный блок}
\addcontentsline{toc}{section}{Тренировочный блок}

\textbf{Уровень: 3 средних, 6 выше среднего, 1 сложная задача.}

\begin{enumerate}
  \item \textbf{Средняя.}
  Расстояние между городами $A$ и $B$ равно $366\unit{км}$.
  Из города $A$ со скоростью $72\unit{км/ч}$ выехал автомобиль.
  Через $30\unit{мин}$ навстречу ему из города $B$ выехал второй
  автомобиль со скоростью $60\unit{км/ч}$.
  На каком расстоянии от города $A$ произойдёт встреча?

  \item \textbf{Средняя.}
  Два пешехода одновременно вышли из одной точки в одном направлении.
  Скорость первого на $1{,}2\unit{км/ч}$ больше скорости второго.
  Через сколько минут расстояние между ними станет равным $420\unit{м}$?

  \item \textbf{Средняя.}
  Длина круговой трассы равна $18\unit{км}$.
  Скорость первого автомобиля равна $48\unit{км/ч}$.
  Через $24\unit{мин}$ после старта второй автомобиль впервые
  обогнал первый.
  Найдите скорость второго автомобиля.

  \item \textbf{Выше среднего.}
  Собственная скорость катера равна $24\unit{км/ч}$,
  скорость течения --- $4\unit{км/ч}$.
  Катер прошёл некоторое расстояние по течению,
  стоял $2\unit{ч}$ и вернулся обратно.
  Всё путешествие заняло $11\unit{ч}$.
  Найдите путь катера за весь рейс.

  \item \textbf{Выше среднего.}
  Первую треть пути автомобиль ехал со скоростью $24\unit{км/ч}$,
  вторую треть --- $30\unit{км/ч}$,
  последнюю треть --- $40\unit{км/ч}$.
  Найдите среднюю скорость на всём пути.

  \item \textbf{Выше среднего.}
  Автомобиль ехал $2\unit{ч}$ со скоростью $50\unit{км/ч}$
  и затем $3\unit{ч}$ со скоростью $70\unit{км/ч}$.
  Найдите среднюю скорость за всё время движения.

  \item \textbf{Выше среднего.}
  Два поезда длиной $180\unit{м}$ и $120\unit{м}$
  движутся навстречу друг другу со скоростями
  $54\unit{км/ч}$ и $36\unit{км/ч}$.
  За сколько секунд они полностью пройдут друг мимо друга?

  \item \textbf{Выше среднего.}
  Поезд длиной $180\unit{м}$ движется со скоростью $72\unit{км/ч}$.
  За сколько секунд он полностью пройдёт мост длиной $420\unit{м}$?

  \item \textbf{Выше среднего.}
  Велосипедист выехал из посёлка со скоростью $12\unit{км/ч}$.
  Через час вслед за ним выехал мотоциклист со скоростью
  $30\unit{км/ч}$.
  На каком расстоянии от посёлка мотоциклист догонит велосипедиста?

  \item \textbf{Сложная.}
  Катер прошёл $36\unit{км}$ по течению реки,
  сделал остановку на $30\unit{мин}$ и вернулся обратно.
  Скорость течения равна $3\unit{км/ч}$,
  всё путешествие заняло $4\unit{ч}$.
  Найдите собственную скорость катера.
\end{enumerate}

%------------------------------------------------
\clearpage
\section*{Ответы}
\addcontentsline{toc}{section}{Ответы}

\begin{table}[ht]
\centering
\caption{Ответы к тренировочным упражнениям}
\begin{tabularx}{\linewidth}{@{}c l X@{}}
\toprule
№ & Ответ & Ключевая проверка \\
\midrule
1
& $216\unit{км}$
& Время первого автомобиля на $0{,}5\unit{ч}$ больше \\

2
& $21\unit{мин}$
& $0{,}42:1{,}2=0{,}35\unit{ч}$ \\

3
& $93\unit{км/ч}$
& Выигрыш одного круга за $0{,}4\unit{ч}$ \\

4
& $210\unit{км}$
& Расстояние в одну сторону равно $105\unit{км}$ \\

5
& $30\unit{км/ч}$
& Три равных участка пути \\

6
& $62\unit{км/ч}$
& Общий путь $310\unit{км}$, общее время $5\unit{ч}$ \\

7
& $12\unit{с}$
& Относительная скорость равна $25\unit{м/с}$ \\

8
& $30\unit{с}$
& Требуемый путь головы поезда равен $600\unit{м}$ \\

9
& $20\unit{км}$
& Догон произойдёт через $\dfrac{2}{3}\unit{ч}$
после старта мотоциклиста \\

10
& $21\unit{км/ч}$
& Ограничение $v>3$; второй корень условию не удовлетворяет \\
\bottomrule
\end{tabularx}
\end{table}

\vfill

\begin{center}
{\color{darkaccent}\rule{0.72\linewidth}{0.6pt}}\\[0.5em]
\small
Проверяйте единицы, относительную скорость
и точную формулировку искомой величины.
\end{center}

\end{document}